\documentclass[10pt,a4paper]{amsart}
\usepackage[margin=19mm]{geometry}
\usepackage{amsmath,amssymb,mathtools}
\usepackage[scr=rsfs,cal=euler]{mathalfa}
\usepackage{xfrac}
\usepackage[unicode,hidelinks,bookmarks=false]{hyperref}
\newcommand{\ie}{i.e.\ }
\newcommand{\cf}{cf.\ }
\newcommand{\resp}{respectively\ }
\newcommand{\Subsection}{\subsection}
\providecommand{\nobreakdash}{\nobreak\hbox{-}}
\setlength{\emergencystretch}{2em}
\multlinegap0pt
\usepackage{bm}
\usepackage{bbm}
\usepackage{dsfont}
\usepackage{xspace}
\usepackage{cancel}
\usepackage{mathtools}
\usepackage{amsthm}
\makeatletter
\@ifundefined{theorem}{\newtheorem{theorem}{Theorem}}{}
\makeatother
\newcommand{\Aut}{\operatorname{Aut}}
\newcommand{\Hol}{\mathrm{Hol}}
\newcommand{\Hom}{\mathrm{Hom}}
\newcommand{\Z}{\mathbf Z}
\title[Determining skew left braces of size np]{Determining skew left braces of size np}
\author{Teresa Crespo, Daniel Gil-Mu\~noz, Anna Rio, Montserrat Vela}
\date{Source: arXiv:2302.13098v3 (CC BY 4.0)}
\begin{document}
\maketitle

\noindent\emph{Source and licence.} Determining skew left braces of size np. Version arXiv:2302.13098v3 (CC BY 4.0), distributed under the Creative Commons Attribution 4.0 International licence, as recorded in the arXiv licence metadata for this exact version and shown on the abstract page https://arxiv.org/abs/2302.13098v3. Full rights notice: licenses/arxiv-2302.13098-CC-BY-4.0.txt.\par\smallskip

\noindent\textit{Editorial scope.} Counted unit: the Theorem whose source label is \texttt{2Gs} in the article (source0.tex lines 727--753, source label \texttt{2Gs}), delivered together with the article's own complete original proof of it (source0.tex lines 755--934, one \texttt{proof} environment of 180 lines). No mathematics is written, repaired or re-derived: statement and proof are the article's own text line for line. Required context retained uncounted: npgivesn (lines 635--639); F (lines 652--660); sgc1 (lines 671--676); sgc3 (lines 671--676). Every definition, display and label of those lines is retained unchanged. Omitted from this package and not redistributed: 1--634, 640--651, 661--670, 677--726, 754--754, 935--2550. Nothing inside a retained excerpt is omitted or altered: every retained body is delivered exactly as the article's lines are written, so no omission receipt of any kind accompanies this package. Citations: the retained excerpts carry no citation command at all, so no bibliography entry of the article is transcribed and no reference is redistributed. Reference adaptation: none was needed and none was made; every reference inside the delivered unit resolves inside it, so no unresolved reference or citation is produced. Printed slips kept literal: no printed word, symbol, hypothesis or number of the counted statement or of its proof was altered. Layout and class: the article's own class is \texttt{article} with packages including latexsym, amsmath, amssymb, amscd, array, comment, xy, amsthm, amsfonts, enumerate, hyperref, mathtools, fullpage, relsize; the wrapper is the lane's pinned \texttt{amsart} editorial wrapper at 10pt on A4, which restates the theorem-like environments the retained text needs and adds the article's own macro definitions that the retained text needs (\texttt{\textbackslash Aut}, \texttt{\textbackslash Hol}, \texttt{\textbackslash Hom}, \texttt{\textbackslash Z}), with extra wrapper packages (bm, bbm, dsfont, xspace, cancel, mathtools). The delivered numbering is the unit's own. Assets: no retained span contains a figure environment, a graphics-inclusion command, a PDF-page inclusion, a table or a TikZ picture; no image file is copied into this package, the package declares no asset dependency and no image file is redistributed with it. Licence: version arXiv:2302.13098v3 of this article is distributed under the Creative Commons Attribution 4.0 International licence, as recorded in the arXiv licence metadata for this exact version and shown on the abstract page https://arxiv.org/abs/2302.13098v3. A full rights notice is delivered at licenses/arxiv-2302.13098-CC-BY-4.0.txt. Characters: every retained excerpt is pure ASCII, so the delivered engine, XeLaTeX with its default Latin Modern text font, reports no missing character.
	\begin{theorem}\label{npgivesn} Let $p$ and $n$ be as in the assumption. Let $(B,\cdot,\circ)$ be a skew brace of size $np$ with additive group $N=\Z_p\rtimes_{\sigma} E$ and multiplicative group $G$, where $(E,\cdot)$ is a group of order $n$ and $\sigma: E\to \Z_p^*=\Aut(\Z_p)$ is a group homomorphism. If $p$ does not divide the order of $\Aut(E)$, then 
		the quotient skew brace $B/\Z_p$ is a skew brace $(B_n, \cdot,\circ)$ with additive group $(E,\cdot)$ such that $\sigma$ is a skew brace homomorphism $B_n\to \Z_p^*$ and $G\simeq \Z_p\rtimes_{\tau} (B_n,\circ)$ for some $\tau: (B_n,\circ)\to \Z_p^*$ group homomorphism. 
		
		An isomorphism class of skew braces of size $np$ with additive group $\Z_p\rtimes_{\sigma} E$ provides an isomorphism class of skew braces of size $n$ with additive group $E$.
	\end{theorem}

\begin{equation}\label{F}		
\hat F=\Big\{ h_a\coloneqq\Big[\begin{pmatrix}
		m_a\\
		a
		\end{pmatrix},\begin{pmatrix}
		\alpha_a & \gamma_{i_a}\\
		0&\lambda_a
		\end{pmatrix}\Big] \, : \, a \in E\Big\},
\end{equation}

\begin{align}
m_{a\lambda_a(b)}&=m_a+\sigma(a)\alpha_am_b+\sigma(a)\gamma_{i_a}(b),\label{sgc1}\\
\alpha_{a\lambda_a(b)}&=\alpha_a\alpha_b,\\
\gamma_{i_{a\lambda_a(b)}}&=\gamma_{i_a}+\alpha_a\gamma_{i_b}, \label{sgc3} \\
\lambda_{a\lambda_a(b)}&=\lambda_a\lambda_b. \label{sgc4}
\end{align}

\begin{theorem}\label{2Gs}
Let $p$ and $n$ be as in the assumption and such that $p\nmid \theta(n)$. Let $B_n$ be a skew brace of size $n$ with additive group $(E,\cdot)$ and let
	$F=\{(a,\lambda_a): a\in E\}$ the corresponding regular subgroup of $\Hol(E)$.
		
	For every $\sigma\in\Hom(B_n,\Z_p^*)$ and every $\tau\in \Hom(F,\Z_p^*)$,
	$$
	G_{(\sigma,\tau)}=\left\{\Big[\begin{pmatrix}m\\a\end{pmatrix},
	\begin{pmatrix}\sigma(a)^{-1}\tau(a,\lambda_a) & 0\\0&\ \lambda_a\end{pmatrix}\Big]
		\ \colon 
		m\in\Z_p,\  a\in E
		\right\}
	$$
	$$G'_{(\sigma,\tau)}=\left\{\Big[\begin{pmatrix}m\\a\end{pmatrix},
		\begin{pmatrix}\tau(a,\lambda_a) & \gamma_{-\sigma(a^{-1})m}\\0&\lambda_a\end{pmatrix}\Big]
		\ \colon 
		m\in\mathbb Z_p,\  a\in E
		\right\}
	$$
	are regular subgroups of $\Hol(\Z_p\rtimes_{\sigma}E)$
	isomorphic to $\Z_p\rtimes_{\tau} F$.
		
	We have $G_{(1,\tau)}=G'_{(1,\tau)}$ for any $\tau$. For a
	non-trivial $\sigma$,
	$G_{(\sigma,\tau)}$ and $G'_{(\sigma,\tau)}$ are not conjugate in $\Hol(\Z_p\rtimes_{\sigma}E)$. Moreover, $G_{(\sigma,\tau)}$
	and $G'_{(\sigma,\tau)}$ represent the unique conjugacy classes of regular 
	subgroups of $\Hol(\Z_p\rtimes_{\sigma} E)$ among the isomorphism class of $\Z_p\rtimes_{\tau}F$ such that the quotient of the corresponding skew brace of size $np$ by the trivial brace $\Z_p$ is isomorphic to $B_n$.
\end{theorem}

\begin{proof}
Since $\sigma$ gives a group homomorphism $E\to \Z_p^*$, we can consider the semidirect product $N=\Z_p\rtimes_{\sigma}E$.  Let $G$ be a regular subgroup of $\Hol(N)$ and $B$ the corresponding skew brace of size $np$. Let us assume that the quotient structure $B/\Z_p$ corresponds to $F$, so that $G\simeq \Z_p\rtimes_{\tau} F$. Following the description of $G$ from the proof of Theorem \ref{npgivesn}, either $G=P_0 \rtimes \hat F$ or $G=P_{-1}\rtimes \hat F$, for some subgroup $\hat F$ as in \eqref{F}. The action given by $F\simeq \hat F\to \Aut(P_i)\simeq \Aut(\Z_p)\simeq \Z_p^*$ is 
$$
(a,\lambda_a)\to h_a\to \Phi_{h_a} \in\Aut(P_i)
$$
followed by the identification of the element in $\Z_p*$ corresponding to the conjugation in $P_i$ by $h_a$.
We obtain 
 $(a,\lambda_a) \mapsto \sigma(a)\alpha_a\in\Z_p^*$ for $i=0$ and $(a,\lambda_a) \mapsto \alpha_a\in \Z_p^*$ for $i=-1$. We want this to agree with the given $\tau:F\to \Z_p^*.$

Let us see that, up to conjugation in $\Hol(E)$, we have actually $G=P_i \rtimes F_0$, where
$$F_0=\Big\{\Big[\begin{pmatrix}
		0\\
		a
		\end{pmatrix},\begin{pmatrix}
		\alpha_a & 0\\
		0&\lambda_a
		\end{pmatrix}\Big] \, : \, a \in E\Big\}\subseteq \Hol(E),$$
with $\lambda_a$ determined by $F$, $\alpha_a=\sigma(a^{-1})\tau(a,\lambda_a)$ for $i=0$ and $\alpha_a=\tau(a,\lambda_a)$ for $i=-1$.

 The equality \eqref{sgc3} is equivalent to $i_{a\lambda_a(b)}=\alpha_a i_b+i_a$, which gives that $a \to i_a$ is a 1-cocycle from $F$ to $\Z_p$ with respect to the action given by $\alpha\colon a\to\alpha_a$, and therefore
 a coboundary. We have then $i_a=j(1-\alpha_a)$ for some fixed $j$. Let us consider first a subgroup $P_0\rtimes \hat F$ with $\hat F$ as in \eqref{F}. We have
		$$\begin{array}{l}
		\Big[\begin{pmatrix}
		0\\
		1
		\end{pmatrix},\begin{pmatrix}
		1 & \gamma_j\\
		0&1
		\end{pmatrix}\Big]\Big[\begin{pmatrix}
		0\\
		a
		\end{pmatrix},\begin{pmatrix}
		\alpha_a & 0\\
		0&\lambda_a
		\end{pmatrix}\Big]
\Big[\begin{pmatrix}
		0\\
		1
		\end{pmatrix},\begin{pmatrix}
		1 & -\gamma_j\\
		0&1
		\end{pmatrix}\Big]\\[12pt]
=\Big[\begin{pmatrix}
		\gamma_j(a)\\
		a
		\end{pmatrix},\begin{pmatrix}
		\alpha_a & -\alpha_a\gamma_j+\gamma_j\\
		0&\lambda_a
		\end{pmatrix}\Big]
=\Big[\begin{pmatrix}
		\gamma_j(a)\\
		a
		\end{pmatrix},\begin{pmatrix}
		\alpha_a & \gamma_{i_a}\\
		0&\lambda_a
		\end{pmatrix}\Big],
\end{array}
		$$
using $-\alpha_a\gamma_j+\gamma_j=(1-\alpha_a)\gamma_j=\gamma_{i_a}$.

On the other hand, inside $P_0\rtimes \hat F$, the subgroup $F$ is determined up to conjugation by $P_0$. Since both $m_a$ and $\gamma_j(a)$ satisfy \eqref{sgc1}, we obtain that $a\mapsto m_a-\gamma_j(a)$ is a 1-cocycle with respect to the action of $F$ on $\Z_p$ given by $\sigma(a)\alpha_a$ and therefore a coboundary:  $m_a-\gamma_j(a)=k(1-\sigma(a)\alpha_a)$ for some fixed integer $k$. By conjugation inside $P_0\rtimes \hat F$, we obtain
		$$\begin{array}{l}
		\Big[\begin{pmatrix}
		k\\
		1
		\end{pmatrix},\begin{pmatrix}
		1 & 0\\
		0&1
		\end{pmatrix}\Big]\Big[\begin{pmatrix}
		\gamma_j(a)\\
		a
		\end{pmatrix},\begin{pmatrix}
		\alpha_a & \gamma_{i_a}\\
		0&\lambda_a
		\end{pmatrix}\Big]
\Big[\begin{pmatrix}
		-k\\
		1
		\end{pmatrix},\begin{pmatrix}
		1 & 0\\
		0&1
		\end{pmatrix}\Big]\\[12pt]
=\Big[\begin{pmatrix}
		k+\gamma_j(a)-\sigma(a)\alpha_a k\\
		a
		\end{pmatrix},\begin{pmatrix}
	\alpha_a & \gamma_{i_a}\\
		0&\lambda_a
		\end{pmatrix}\Big]
=\Big[\begin{pmatrix}
		m_a\\
		a
		\end{pmatrix},\begin{pmatrix}
	\alpha_a & \gamma_{i_a}\\
		0&\lambda_a
		\end{pmatrix}\Big].
        \end{array}
		$$
To sum up, we have obtained that, given any $\hat F$ as in \eqref{F},  $P_0\rtimes F_0$ is conjugate to $P_0\rtimes \hat F$ inside $\Hol(N)$.

Let us consider now  $P_{-1}\rtimes \hat F$ with $\hat F$ as in \eqref{F}. For any integer $k$, we have
		$$\begin{array}{l}
		\Big[\begin{pmatrix}
		0\\
		1
		\end{pmatrix},\begin{pmatrix}
		1 & \gamma_{j+k}\\
		0&1
		\end{pmatrix}\Big]\Big[\begin{pmatrix}
		0\\
		a
		\end{pmatrix},\begin{pmatrix}
		\alpha_a & 0\\
		0&\lambda_a
		\end{pmatrix}\Big]
\Big[\begin{pmatrix}
		0\\
		1
		\end{pmatrix},\begin{pmatrix}
		1 & -\gamma_{j+k}\\
		0&1
		\end{pmatrix}\Big]\\[12pt]
=\Big[\begin{pmatrix}
		\gamma_{j+k}(a)\\
		a
		\end{pmatrix},\begin{pmatrix}
		\alpha_a & -\alpha_a\gamma_{j+k}+\gamma_{j+k}\\
		0&\lambda_a
		\end{pmatrix}\Big]
=\Big[\begin{pmatrix}
		\gamma_{j+k}(a)\\
		a
		\end{pmatrix},\begin{pmatrix}
		\alpha_a & \gamma_{i_a}+(1-\alpha_a)\gamma_k\\
		0&\lambda_a
		\end{pmatrix}\Big].
\end{array}
		$$
As above, $m_a-\gamma_j(a)=k(1-\sigma(a)\alpha_a)$ for some fixed integer $k$. By conjugation inside $P_{-1} \rtimes \hat F$, we obtain
		$$\begin{array}{l}
		\Big[\begin{pmatrix}
		k\\
		1
		\end{pmatrix},\begin{pmatrix}
		1 & \gamma_{-k}\\
		0&1
		\end{pmatrix}\Big]\Big[\begin{pmatrix}
		\gamma_{j+k}(a)\\
		a
		\end{pmatrix},\begin{pmatrix}
		\alpha_a & \gamma_{i_a}+(1-\alpha_a)\gamma_k\\
		0&\lambda_a
		\end{pmatrix}\Big]
\Big[\begin{pmatrix}
		-k\\
		1
		\end{pmatrix},\begin{pmatrix}
		1 & \gamma_k\\
		0&1
		\end{pmatrix}\Big]\\[12pt]
=\Big[\begin{pmatrix}
		k+\gamma_j(a)-\sigma(a)\alpha_a k\\
		a
		\end{pmatrix},\begin{pmatrix}
	\alpha_a & \gamma_{i_a}\\
		0&\lambda_a
		\end{pmatrix}\Big]
=\Big[\begin{pmatrix}
		m_a\\
		a
		\end{pmatrix},\begin{pmatrix}
	\alpha_a & \gamma_{i_a}\\
		0&\lambda_a
		\end{pmatrix}\Big].
        \end{array}
		$$
and we see that, given any $\hat F$ as in \eqref{F},  $P_{-1}\rtimes F_0$ is conjugate to $P_{-1}\rtimes \hat F$ inside $\Hol(N)$.

In conclusion, we have seen that $G_{(\sigma,\tau)}=P_0\rtimes F_0$ and $G'_{(\sigma,\tau)}=P_{-1}\rtimes F_0$ are, up to conjugation, the unique regular subgroups of $\Hol(N)$ giving $B_n$ by the quotient procedure.
\end{proof}

\end{document}
